\chapterimage{chapter_head_1.pdf}
\chapter{Exam Prep \& Programming Assignments}

\section{High-Frequency Exam Topics}

Success in CS1 and the Foundation Exam requires mastery of both theoretical analysis and applied data structure operations.

\begin{theorem}[Algorithmic Analysis]
    \begin{itemize}
        \item \textbf{Recurrence Relations:} Mastery of the Master Theorem and the Iteration Technique is mandatory.
        \item \textbf{Summations:} Algebraic manipulation of closed-form solutions.
        \item \textbf{Runtime Analysis:} Determining Big-O for nested loops and recursive calls.
    \end{itemize}
\end{theorem}

\begin{definition}[Core Data Structure Operations]
    \begin{itemize}
        \item \textbf{Binary Heaps:} Tracing insertions and deletions in an array-based heap.
        \item \textbf{BST \& Tries:} Implementing recursion for counting or searching.
        \item \textbf{Linked Lists:} Pointer manipulation for insertion, deletion, and reordering.
    \end{itemize}
\end{definition}

\section{Programming Assignment Core Logic}

Each Programming Assignment (PA) introduces a specific algorithmic "trick" that builds on standard data structures.

\begin{example}[PA Hybrid Sorting]
    In PA4, Merge Sort and Quick Sort were optimized by switching to \textbf{Insertion Sort} once the subarray size fell below a certain threshold (e.g., 30). This leverages the efficiency of $O(n^2)$ algorithms on small datasets.
\end{example}

\begin{example}[PA Order Statistic Tree]
    In PA5, a standard BST was augmented with a \texttt{subtree\_size} field at each node. This allowed for $O(h)$ time complexity when searching for the \textbf{$k$-th smallest element}.
\end{example}

\begin{example}[PA Mode-Dependent Heap]
    In PA6, a Binary Heap was implemented with a mode-dependent comparator, allowing it to toggle between a \textbf{Min-Heap} (Triage Mode) and a \textbf{Max-Heap} (Adoption Mode) using a bottom-up $O(n)$ \texttt{heapify}.
\end{example}

\section{Foundation Exam Core Patterns}

\begin{remark}[The Recursive Traversal Pattern]
    Most tree and graph problems follow a standard template:
    \begin{enumerate}
        \item \textbf{Base Case:} Check for \texttt{NULL} or terminal states.
        \item \textbf{Recursive Step:} Aggregate results from child nodes.
        \item \textbf{Return:} Pass the result back up the stack.
    \end{enumerate}
\end{remark}

\begin{remark}[Pointer Safety]
    In Linked List and DMA questions, always trace your pointers. Ensure you don't "drop" the rest of the list when reassigning \texttt{next} pointers.
\end{remark}

\section{Practice Problems}

\begin{exercise}[Recurrence Iteration]
    Solve the recurrence $T(n) = 2T(n/2) + n$ using the iteration technique. Show each step of the expansion.
\end{exercise}

\begin{exercise}[Heap Trace]
    Given an empty min-heap, trace the insertion of the following values: \texttt{[10, 5, 12, 3, 8]}. Draw the resulting array representation.
\end{exercise}
